\section{Results}
\label{sec:results}

\Cref{tab:main} gives the headline numbers and \Cref{tab:bytype} the per-category
breakdown. \Cref{fig:bytype} plots the latter and \Cref{tab:grid} shows every one
of the 180 attempts individually, unaggregated.

Summed over the question set, GraphRAG solved 14 of 20 questions on all three
attempts and passed 45 of 60 attempts; text-to-SQL solved 9 and passed 29; vector
RAG solved 3 and passed 9. These totals weight four hand-chosen categories
equally and GraphRAG lost two of the four, so the aggregate is a property of this
question mix rather than an overall ranking of the architectures. The per-category
results are the informative ones.

\subsection{By category}

\textbf{T1, single-document semantic.} GraphRAG passed 15 of 15. Vector RAG, the
expected winner, passed 3 of 15 --- all three from a single question --- with 6
refusals and 6 answers that named the right company and the right risks but cited
no filing. Text-to-SQL passed 0 of 15 despite reaching the correct text in 14
attempts: its generated queries used full-text ranking to find the passage but did
not select the accession number, so the answers were right and unciteable. This is
the clearest demonstration that provenance is a separate capability from
retrieval.

\textbf{T2, structured aggregation.} Text-to-SQL passed 15 of 15, as predicted.
GraphRAG passed 6 of 15, solving 2 of 5 questions, with 7 attempts returning wrong
answers and 2 refusing. Vector RAG passed 0 of 15 and refused all fifteen times,
which is the correct behaviour: an aggregate over 10{,}511 companies is not
present in any eight retrieved chunks, and the arm said so rather than guessing.

\textbf{T3, multi-hop relational plus text.} GraphRAG passed 15 of 15; both
baselines passed 0 of 15. This is the category the design targets and the result
is unambiguous. Text-to-SQL produced the benchmark's only two falsehoods here.

\textbf{T4, cross-document entity resolution.} Written for the graph, and the
graph came second. Text-to-SQL passed 14 of 15 and solved 4 of 5 questions;
GraphRAG passed 9 of 15 and solved 2 of 5 with 2 more solved inconsistently, its
6 misses all refusals rather than errors; vector RAG passed 6 of 15.

\subsection{Outcome composition}

\Cref{fig:outcomes} decomposes all 60 attempts per arm. The composition matters as
much as the pass rate, because the failure modes are not interchangeable.

Vector RAG refused 45 times. That is a low score but a benign failure: asked about
a company whose chunks it could not find, it generally said the excerpts did not
contain the answer rather than answering about whichever company did surface.
Text-to-SQL was uncited 16 times --- right, and unverifiable. GraphRAG was never
uncited, which follows from its construction: the traversal returns accession
numbers before any text is read, so every claim arrives with a key attached.

Only \textbf{2 of the 180 attempts} asserted something verifiably false, both
text-to-SQL on T3.

\subsection{Cost and latency}

\Cref{fig:cost} plots tokens against attempts passed. GraphRAG consumed
780{,}500 tokens to text-to-SQL's 158{,}171 and vector RAG's 118{,}850 --- about
17{,}300 tokens per passing attempt against 5{,}450 and 13{,}200 respectively. Its
accuracy is bought, not free. The extreme case is T2, where GraphRAG spent
590{,}453 tokens against text-to-SQL's 24{,}604: roughly 24 times the cost for
40\% of the accuracy, on the category it was expected to lose.

The reverse cost is larger and easier to overlook. \Cref{fig:latency} plots mean
latency on a log scale. Text-to-SQL averaged 25.3\,s overall against GraphRAG's
5.9\,s and vector RAG's 1.6\,s, but the overall mean hides the shape: on T3 it
averaged 92\,s and its worst single attempt took 445\,s. The cause is not the
model but the SQL --- the generated queries join \texttt{subsidiary},
\texttt{filing\_section} and \texttt{reporting\_owner} without narrowing first.
Unbounded, LLM-authored SQL is a cost risk against one's own warehouse
independently of whether the answer is correct.

\subsection{Repair usage}

\Cref{tab:repair} and \Cref{fig:repair} report how often each arm used its single
repair allowance. GraphRAG triggered a repair on 34 of 60 attempts against
text-to-SQL's 7 --- nearly five times as often --- and on T4 it needed one on
\emph{all fifteen} attempts. Vector RAG has no query to repair and used none.

This asymmetry accounts for much of the token gap, and it should temper the
reading of GraphRAG's wins: writing correct Cypher against this schema was harder
for the model than writing correct SQL, and the graph arm leaned on its second
chance far more heavily. As noted in \Cref{sec:design}, the zero-row trigger is
also not neutral between the two languages, and the direction of that bias favours
the graph.
